\documentclass{amsart}
\usepackage{amsmath,amssymb,amsthm,hyperref}

\newtheorem{theorem}{Theorem}
\newtheorem{proposition}{Proposition}
\newtheorem{lemma}{Lemma}
\newtheorem{corollary}{Corollary}
\theoremstyle{definition}
\newtheorem{remark}{Remark}
\newtheorem{assumption}{Assumption}

\newcommand{\R}{\mathbb{R}}
\newcommand{\dd}{\,\mathrm{d}}
\DeclareMathOperator{\supp}{supp}

\begin{document}

\title{Nonlinear eigenvalue problems for a local epigenetic structured population equation}
\maketitle

\section*{Setup and notation}

We study the structured population (growth--fragmentation type) equation
\begin{equation}\label{eq:master}
\partial_t f(t,x) = -\bigl(\beta(x)+\mu(x)\bigr) f(t,x) + 2\int_{\R} \beta(y)\,p(x,y)\, f(t,y)\dd y,
\qquad x\in\R,
\end{equation}
with local transition kernel $p(x,y)=\varphi(y-x)$, where $\varphi$ is a probability
density supported in $(-\varepsilon,\varepsilon)$ with
\begin{equation}\label{eq:moments}
\int_\R \varphi(r)\dd r=1,\qquad
a=\int_\R r\,\varphi(r)\dd r,\qquad
D=\int_\R r^2\,\varphi(r)\dd r .
\end{equation}
Throughout, $\beta,\mu\in C^\infty(\R)$ with $\beta\ge 0$, $\mu\ge 0$. We write the
\emph{net reaction coefficient}
\begin{equation}\label{eq:m}
m(x):=\beta(x)-\mu(x).
\end{equation}

\begin{assumption}\label{ass:confine}
$\beta,\mu\in C^\infty(\R)$, $\beta\ge0$, $\mu\ge0$, $D>0$, and $m=\beta-\mu$ is bounded above with
\begin{equation}\label{eq:confine}
m(x)\to-\infty\quad\text{as }|x|\to\infty .
\end{equation}
\end{assumption}

Condition \eqref{eq:confine} is the standard ``confining'' condition guaranteeing a
discrete spectrum and an integrable principal eigenfunction; it is what makes parts
(b)--(e) well posed on the whole line. (If instead one works on a bounded interval
$\Omega$ with Dirichlet data, Assumption~\ref{ass:confine} can be dropped and all
results below hold verbatim with $\R$ replaced by $\Omega$ and the decay conditions
replaced by $F|_{\partial\Omega}=0$; see Remark~\ref{rem:bounded}.)

\bigskip

\section{(a) Local (diffusion) approximation of the integral term}

\begin{proposition}\label{prop:expansion}
Let $\beta\in C^2(\R)$ and let $f(t,\cdot)\in C^2(\R)$ for each $t$. With $a,D$ as in
\eqref{eq:moments}, the integral term in \eqref{eq:master} satisfies, pointwise in $x$,
\begin{equation}\label{eq:Iexp}
2\int_\R \beta(y)\,\varphi(y-x)\,f(t,y)\dd y
= 2\,\beta(x)f(t,x) + 2a\,\partial_x\!\bigl(\beta f\bigr)(t,x)
+ D\,\partial_x^2\!\bigl(\beta f\bigr)(t,x) + R_\varepsilon(t,x),
\end{equation}
where
\begin{equation}\label{eq:remainder}
|R_\varepsilon(t,x)| \le \tfrac{2}{3}\,\varepsilon^3\,
\sup_{|r|\le\varepsilon}\bigl|\partial_x^3(\beta f)(t,x+r)\bigr|
= O(\varepsilon^3).
\end{equation}
Consequently \eqref{eq:master} is, to second order in $\varepsilon$,
\begin{equation}\label{eq:pde-cons}
\boxed{\;\partial_t f = \bigl(\beta-\mu\bigr) f + 2a\,\partial_x(\beta f) + D\,\partial_x^2(\beta f)\;}
\end{equation}
a drift--diffusion equation in conservation (Fokker--Planck) form, with reaction
coefficient $m=\beta-\mu$, drift carried by $2a$ and diffusion by $D$.
\end{proposition}

\begin{proof}
Substitute $r=y-x$ (so $y=x+r$, $\dd y=\dd r$) and write $g(t,\cdot):=\beta(\cdot)f(t,\cdot)$.
Since $\supp\varphi\subset(-\varepsilon,\varepsilon)$,
\[
2\int_\R \beta(y)\varphi(y-x)f(t,y)\dd y
= 2\int_{-\varepsilon}^{\varepsilon} g(t,x+r)\,\varphi(r)\dd r .
\]
By Taylor's theorem with integral remainder, for $|r|<\varepsilon$,
\[
g(t,x+r)= g(t,x) + r\,\partial_x g(t,x) + \tfrac12 r^2\,\partial_x^2 g(t,x)
+ \tfrac12\int_0^r (r-s)^2\,\partial_x^3 g(t,x+s)\dd s .
\]
Multiply by $2\varphi(r)$ and integrate over $(-\varepsilon,\varepsilon)$. Using
\eqref{eq:moments} ($\int\varphi=1$, $\int r\varphi=a$, $\int r^2\varphi=D$), the first three
terms give exactly $2g + 2a\,\partial_x g + D\,\partial_x^2 g$, i.e.\ the explicit right-hand
side of \eqref{eq:Iexp}. The remainder is
\[
R_\varepsilon(t,x)= \int_{-\varepsilon}^{\varepsilon}\!\varphi(r)\!\int_0^r (r-s)^2\,\partial_x^3 g(t,x+s)\dd s\,\dd r .
\]
For $|r|\le\varepsilon$ the inner integral is bounded by
$\frac{|r|^3}{3}\sup_{|s|\le\varepsilon}|\partial_x^3 g(t,x+s)|\le
\frac{\varepsilon^3}{3}\sup_{|s|\le\varepsilon}|\partial_x^3 g(t,x+s)|$. Since
$\int_{-\varepsilon}^\varepsilon \varphi=1$, this gives \eqref{eq:remainder}. Inserting
\eqref{eq:Iexp} into \eqref{eq:master} and combining $-\beta f + 2\beta f=\beta f$ with
$-\mu f$ gives $m f + 2a\,\partial_x(\beta f)+D\,\partial_x^2(\beta f)$, which is
\eqref{eq:pde-cons} up to $R_\varepsilon=O(\varepsilon^3)$.
\end{proof}

\begin{remark}[Relation to the form quoted in part (b)]\label{rem:form}
Equation \eqref{eq:pde-cons} is the \emph{exact} second-order local approximation in
\emph{conservation form}: diffusion coefficient $D$, drift $2a$, with the state dependence
of $\beta$ kept inside the derivatives. The reduced \emph{non-conservative,
frozen-coefficient} form quoted in the problem,
\begin{equation}\label{eq:reduced}
\tfrac{D}{2}\,\partial_x^2 f - a\,\partial_x f + \bigl(\beta-\mu\bigr) f,
\end{equation}
is obtained from \eqref{eq:pde-cons} in the standard \emph{local-transition rescaling}, in
which $\beta$ varies slowly on the transition scale $\varepsilon$ (so
$\partial_x(\beta f)\approx \beta\,\partial_x f$ and $\partial_x^2(\beta f)\approx
\beta\,\partial_x^2 f$, the slowly-varying $\beta$ absorbed into a normalisation of the local
diffusivity), together with the Fokker--Planck diffusion convention $\widetilde D:=D/2$ (half
the second moment) and drift normalisation $\widetilde a:=a$. We therefore analyse the
\emph{general operator}
\begin{equation}\label{eq:Ldef}
\mathcal{A}F := \frac{\widetilde D}{2}\,F'' - \widetilde a\,F' + m\,F,
\qquad \widetilde D>0,\ \widetilde a\in\R,\ m=\beta-\mu,
\end{equation}
of which \eqref{eq:reduced} (with $\widetilde D=D$, $\widetilde a=a$) is the special case.
All qualitative conclusions in (b)--(e) are independent of the particular positive value of
$\widetilde D$ and of $\widetilde a\in\R$.
\end{remark}

\section{(b) The separated steady state and its ODE}

\begin{proposition}\label{prop:ode}
Suppose $f(t,x)=e^{\lambda t}F(x)$ with $F\in C^2(\R)$, $F\ge0$, $\int_\R F=1$,
$F(x)\to0$ as $|x|\to\infty$, solves the local model in the reduced regime of
Remark~\ref{rem:form}. Then $(\lambda,F)$ satisfies the eigenvalue ODE
\begin{equation}\label{eq:ode}
\frac{\widetilde D}{2}\,F''(x) - \widetilde a\,F'(x) + \bigl(\beta(x)-\mu(x)\bigr)F(x) = \lambda\,F(x),
\qquad x\in\R,
\end{equation}
with $\int_\R F=1$ and $F(x)\to0$ as $|x|\to\infty$. In the original moments this reads
$\frac{D}{2}F'' - aF' + (2\beta-\beta-\mu-\lambda)F=0$, since $2\beta-\beta=\beta$, which is
exactly the form stated in the problem.
\end{proposition}

\begin{proof}
Insert $f(t,x)=e^{\lambda t}F(x)$ into
$\partial_t f = \frac{\widetilde D}{2} f'' - \widetilde a f' + m f$ (Remark~\ref{rem:form}).
The left side is $\lambda e^{\lambda t}F$; the right side is
$e^{\lambda t}(\frac{\widetilde D}{2}F'' - \widetilde a F' + m F)$. Cancelling
$e^{\lambda t}\neq0$ gives \eqref{eq:ode}. The constraints $\int F=1$, $F\to0$ are inherited
from the ansatz. Writing $m=2\beta-\beta-\mu$ and moving $\lambda F$ left gives the stated form.
\end{proof}

\section{(c) The principal eigenpair: existence, uniqueness, positivity, regularity, decay}
\label{sec:c}

We analyse $\mathcal{A}F=\lambda F$ with $\mathcal A$ from \eqref{eq:Ldef}, under
Assumption~\ref{ass:confine}.

\subsection*{Self-adjoint reduction (Liouville transform)}
Set
\begin{equation}\label{eq:gauge}
\theta:=\frac{\widetilde a}{\widetilde D},\qquad
F(x)=e^{\theta x}\,u(x),\qquad u(x):=e^{-\theta x}F(x).
\end{equation}
Then $F'=e^{\theta x}(u'+\theta u)$ and $F''=e^{\theta x}(u''+2\theta u'+\theta^2 u)$, so
\[
\frac{\widetilde D}{2}F''-\widetilde a F'+mF
= e^{\theta x}\!\left[\frac{\widetilde D}{2}u'' + \bigl(\widetilde D\theta-\widetilde a\bigr)u'
+\Bigl(\tfrac{\widetilde D}{2}\theta^2-\widetilde a\theta+m\Bigr)u\right].
\]
Since $\widetilde D\theta-\widetilde a=0$ and
$\frac{\widetilde D}{2}\theta^2-\widetilde a\theta=-\frac{\widetilde a^2}{2\widetilde D}$,
the problem $\mathcal A F=\lambda F$ becomes the \emph{self-adjoint Schr\"odinger eigenvalue problem}
\begin{equation}\label{eq:schrod}
\mathcal L u:=\frac{\widetilde D}{2}\,u'' + V(x)\,u = \lambda u,
\qquad
V(x):=m(x)-\frac{\widetilde a^2}{2\widetilde D}=\beta(x)-\mu(x)-\frac{\widetilde a^2}{2\widetilde D}.
\end{equation}
By Assumption~\ref{ass:confine}, $V$ is smooth, bounded above, and $V(x)\to-\infty$ as
$|x|\to\infty$; equivalently $H:=-\mathcal L=-\frac{\widetilde D}{2}\partial_x^2-V$ is a
Schr\"odinger operator with smooth, bounded-below confining potential $-V\to+\infty$.

\subsection*{Functional-analytic framework}
Let $H=-\mathcal L$ act on $L^2(\R)$ with form domain
$Q(H)=\{u\in H^1(\R): \int (-V)\,u^2<\infty\}$, a Hilbert space since $-V$ is bounded below.
The form
\begin{equation}\label{eq:form}
q(u)=\int_\R\Bigl(\frac{\widetilde D}{2}|u'|^2 - V u^2\Bigr)\dd x
\end{equation}
is closed and semibounded below ($\sup V<\infty$), and $H$ is its self-adjoint Friedrichs
operator (Reed--Simon, \emph{Methods of Modern Mathematical Physics} II, Thm.~X.23; IV,
Thm.~XIII.16).

\begin{lemma}[Discreteness and bottom eigenvalue]\label{lem:discrete}
Under Assumption~\ref{ass:confine}, $H=-\mathcal L$ has compact resolvent; hence its spectrum
is purely discrete, consisting of simple\footnote{Simplicity of the $L^2$ eigenvalues of a
one-dimensional Schr\"odinger operator is classical: two linearly independent solutions of a
second-order ODE cannot both be $L^2$ at $+\infty$ (Berezin--Shubin, \emph{The Schr\"odinger
Equation}, \S2).} eigenvalues $E_0<E_1<E_2<\cdots\to+\infty$. The lowest eigenvalue is
\begin{equation}\label{eq:rayleigh}
E_0=\inf_{0\neq u\in Q(H)}\frac{q(u)}{\|u\|_{2}^2},\qquad
\lambda:=-E_0=\sup_{0\neq u\in Q(H)}\frac{\displaystyle\int_\R\!\Bigl(V u^2-\frac{\widetilde D}{2}|u'|^2\Bigr)}{\displaystyle\int_\R u^2}.
\end{equation}
\end{lemma}

\begin{proof}
For a confining potential $-V\to+\infty$, the embedding $Q(H)\hookrightarrow L^2(\R)$ is
compact (Reed--Simon IV, Thm.~XIII.67), so $H$ has compact resolvent; a self-adjoint,
semibounded operator with compact resolvent has purely discrete spectrum accumulating only at
$+\infty$, and the min--max principle gives \eqref{eq:rayleigh} for the bottom eigenvalue.
Since $\mathcal L u=\lambda u\iff Hu=-\lambda u$, the largest eigenvalue $\lambda=-E_0$ of
$\mathcal L$ corresponds to the bottom $E_0$ of $H$, yielding the second formula.
\end{proof}

\begin{theorem}[Principal eigenpair]\label{thm:principal}
Under Assumption~\ref{ass:confine}, problem \eqref{eq:ode} (equivalently \eqref{eq:schrod})
has a \emph{principal eigenpair} $(\lambda,F)$ with:
\begin{enumerate}
\item[\rm(i)] \textbf{(Existence \& value)} $\lambda=-E_0\in\R$ is the largest eigenvalue, given by \eqref{eq:rayleigh}.
\item[\rm(ii)] \textbf{(Positivity \& uniqueness)} The principal eigenfunction $u_0$ of
\eqref{eq:schrod} can be taken $>0$ everywhere; $\lambda$ is simple, so $u_0$ and hence
$F=e^{\theta x}u_0>0$ ($\theta=\widetilde a/\widetilde D$) is the \emph{unique} constant-sign
eigenfunction up to a positive scalar. The normalisation $\int_\R F=1$ fixes $F>0$ uniquely.
\item[\rm(iii)] \textbf{(Regularity)} $F\in C^\infty(\R)$.
\item[\rm(iv)] \textbf{(Decay)} $u_0$ (hence $F$, up to the fixed factor $e^{\theta x}$) decays
faster than any exponential; in particular $|F(x)|\le C_N(1+|x|)^{-N}$ for every $N$, and all
derivatives decay likewise.
\end{enumerate}
\end{theorem}

\begin{proof}
\emph{(i)} is Lemma~\ref{lem:discrete}.

\emph{(ii) Positivity.} The ground state of a one-dimensional confining Schr\"odinger operator
has no zeros: by the Sturm oscillation theorem the $n$-th eigenfunction has exactly $n$ zeros
(Berezin--Shubin \S2.4), so the bottom eigenfunction $u_0$ ($n=0$) has none; take $u_0>0$.
Uniqueness up to a scalar is the simplicity of $\lambda$ (Lemma~\ref{lem:discrete}). Then
$F=e^{\theta x}u_0>0$, and $\int_\R F=1$ fixes the positive scalar uniquely.
\emph{(Equivalently:} for $c$ large, $(H+c)^{-1}$ is positivity-improving by the maximum
principle, so by Krein--Rutman its spectral radius is a simple eigenvalue with a positive
eigenfunction, namely $u_0$.)

\emph{(iii) Regularity.} $V\in C^\infty(\R)$. From $\frac{\widetilde D}{2}u_0''=(\lambda-V)u_0$,
if $u_0\in C^k$ then $u_0''=\frac{2}{\widetilde D}(\lambda-V)u_0\in C^k$, so $u_0\in C^{k+2}$.
Starting from $u_0\in H^1\subset C^{0,1/2}_{loc}$ (1D Sobolev) and iterating, $u_0\in C^\infty(\R)$,
hence $F=e^{\theta x}u_0\in C^\infty(\R)$.

\emph{(iv) Decay.} Fix $\lambda$. By \eqref{eq:confine}, $V(x)-\lambda\to-\infty$, so for any
$A>0$ there is $R_A$ with $V(x)-\lambda\le -A$ for $|x|\ge R_A$. There the equation
$u_0''=\frac{2}{\widetilde D}(\lambda-V)u_0\ge\frac{2A}{\widetilde D}u_0>0$ holds with $u_0>0$.
The comparison function $w_A(x)=u_0(R_A)e^{-\sqrt{2A/\widetilde D}\,(|x|-R_A)}$ satisfies
$w_A''=\frac{2A}{\widetilde D}w_A$, agrees with $u_0$ at $x=R_A$, and dominates $u_0$ at
$+\infty$ (both $\to0$); the comparison/maximum principle for $u_0''\ge\frac{2A}{\widetilde D}u_0$
gives $0<u_0(x)\le w_A(x)$ for $|x|\ge R_A$. As $A>0$ is arbitrary,
$u_0(x)=O(e^{-c|x|})$ for \emph{every} $c>0$, so $|u_0(x)|\le C_N(1+|x|)^{-N}$ for all $N$.
The fixed factor $e^{\theta x}$ in $F=e^{\theta x}u_0$ does not affect this faster-than-exponential
decay. Differentiating the equation and bootstrapping gives the same decay for all
derivatives, and $\int_\R F<\infty$ (consistency of $\int F=1$).
\end{proof}

\begin{remark}[The nonlinear eigenvalue problem and the conservation form]
For the full conservation form \eqref{eq:pde-cons}, the steady equation
$\lambda F=mF+2a(\beta F)'+D(\beta F)''$ becomes, for the unknown $G:=\beta F$ ($\beta>0$),
$D G''+2aG'+mG=\lambda G/\beta\cdot\beta$, i.e.\ a problem of the form \eqref{eq:Ldef} with
$\widetilde D=2D$, $\widetilde a=-2a$, potential $m$; Theorem~\ref{thm:principal} applies at
fixed coefficients. The map $(\beta,\mu)\mapsto(\lambda,F)$ is nonlinear (the coefficients
appear nonlinearly), but for each fixed smooth $(\beta,\mu)$ the spectral problem is linear,
which is what we analyse.
\end{remark}

\subsection*{Dependence of $\lambda$ on $\beta,\mu,a,D$}

\begin{proposition}[Monotonicity and dependence]\label{prop:dependence}
Let $\lambda=\lambda[m,\widetilde a,\widetilde D]$ be the principal eigenvalue. Then:
\begin{enumerate}
\item[\rm(M1)] \textbf{(Monotone in $m$)} $\lambda$ is nondecreasing in $m=\beta-\mu$: $m_1\le m_2$
pointwise $\Rightarrow$ $\lambda[m_1]\le\lambda[m_2]$, strictly if $m_1<m_2$ on a positive-measure
set. Thus $\lambda$ increases with $\beta$ and decreases with $\mu$.
\item[\rm(M2)] \textbf{(Drift)} $\lambda[\widetilde a]=\lambda[0]-\widetilde a^2/(2\widetilde D)$;
hence $\lambda$ is maximal at $\widetilde a=0$ and strictly decreasing in $|\widetilde a|$.
\item[\rm(M3)] \textbf{(Diffusion)} At $\widetilde a=0$, $\widetilde D\mapsto\lambda$ is concave and
strictly decreasing, with
$\partial_{\widetilde D}\lambda=-\tfrac12\int|u_0'|^2/\!\int u_0^2\le0$.
\item[\rm(M4)] \textbf{(Bounds)} $\displaystyle\int_\R m\,F\le\sup_\R m$ and
$\lambda\le\sup_\R m-\widetilde a^2/(2\widetilde D)$; a Gaussian trial function centred at
$\arg\max m$ gives a matching lower bound up to the local curvature of $m$.
\end{enumerate}
\end{proposition}

\begin{proof}
By \eqref{eq:rayleigh}, with $\widetilde a$ absorbed into $V=m-\widetilde a^2/(2\widetilde D)$,
\begin{equation}\label{eq:lamvar}
\lambda=\sup_{\|u\|_2=1}\ \int_\R\Bigl(m\,u^2-\tfrac{\widetilde D}{2}|u'|^2\Bigr)\dd x
-\frac{\widetilde a^2}{2\widetilde D}.
\end{equation}

(M1) For fixed $u$, $u\mapsto\int(mu^2-\frac{\widetilde D}{2}|u'|^2)$ is nondecreasing in $m$
($u^2\ge0$); the supremum preserves this. Strictness: with $u_1>0$ the principal eigenfunction
for $m_1$, $\lambda[m_2]\ge\int(m_2u_1^2-\frac{\widetilde D}{2}|u_1'|^2)
>\int(m_1u_1^2-\frac{\widetilde D}{2}|u_1'|^2)=\lambda[m_1]$, since $\int(m_2-m_1)u_1^2>0$.

(M2) is the additive shift in \eqref{eq:lamvar}: the supremum is over $u$ independent of
$\widetilde a$.

(M3) At $\widetilde a=0$, $\lambda(\widetilde D)$ is a supremum of affine functions of
$\widetilde D$, hence concave. By Feynman--Hellmann (differentiating
$\lambda=\langle u_0,\mathcal L u_0\rangle/\|u_0\|_2^2$ and using that first-order variations
of the optimiser $u_0$ vanish), $\partial_{\widetilde D}\lambda
=\langle u_0,\frac12\partial_x^2u_0\rangle/\|u_0\|_2^2=-\frac12\int|u_0'|^2/\!\int u_0^2\le0$,
strict since $u_0\in L^2$, $u_0>0$ forces $u_0'\not\equiv0$.

(M4) From \eqref{eq:lamvar} with unit optimiser $u_0$,
$\lambda=\int mu_0^2-\frac{\widetilde D}{2}\int|u_0'|^2-\frac{\widetilde a^2}{2\widetilde D}
\le\sup m-\frac{\widetilde a^2}{2\widetilde D}$. Also $\int mF\le\sup m$ since $F\ge0$,
$\int F=1$. The lower bound is $\lambda\ge\int m\psi^2-\frac{\widetilde D}{2}\int|\psi'|^2
-\frac{\widetilde a^2}{2\widetilde D}$ for any unit $\psi$, e.g.\ a narrow Gaussian at
$\arg\max m$, yielding $\sup m-\frac{\widetilde a^2}{2\widetilde D}-O(\widetilde D^{1/2})$.
\end{proof}

\begin{remark}[Bounded-domain version]\label{rem:bounded}
On a bounded interval $\Omega$ with Dirichlet conditions $F|_{\partial\Omega}=0$,
Assumption~\ref{ass:confine} is unnecessary: $\mathcal L$ on $\Omega$ has compact resolvent
(Rellich), and Theorem~\ref{thm:principal} and Proposition~\ref{prop:dependence} hold with
$\R\to\Omega$, decay replaced by $F|_{\partial\Omega}=0$, and $F>0$ in $\Omega$.
\end{remark}

\section{(d) Entropy / Lyapunov laws and long-time convergence}

We use the reduced equation $\partial_t f=\mathcal A f$, $\mathcal A$ from \eqref{eq:Ldef}, and
its adjoint $\mathcal A^*g=\frac{\widetilde D}{2}g''+\widetilde a g'+mg$.

\begin{lemma}[Direct and dual eigenfunctions]\label{lem:dual}
Let $u_0>0$ be the principal eigenfunction of \eqref{eq:schrod}
($\mathcal L u_0=\lambda u_0$) and $\theta=\widetilde a/\widetilde D$. Then
\begin{equation}\label{eq:Fphi}
F(x)=e^{\theta x}u_0(x)>0,\qquad \phi(x)=e^{-\theta x}u_0(x)>0
\end{equation}
satisfy $\mathcal A F=\lambda F$ and $\mathcal A^*\phi=\lambda\phi$, with $\phi$ the principal
eigenfunction of the adjoint. Moreover $\phi F=u_0^2>0$, and after rescaling we may assume
$\int_\R \phi F\,\dd x=\int_\R u_0^2\,\dd x=1$.
\end{lemma}

\begin{proof}
$\mathcal A F=\lambda F$ is the Liouville reduction \S\ref{sec:c}. For the adjoint, substitute
$\phi=e^{-\theta x}v$: a direct computation gives
$\mathcal A^*\phi=e^{-\theta x}\bigl(\frac{\widetilde D}{2}v''+(m-\frac{\widetilde a^2}{2\widetilde D})v\bigr)
=e^{-\theta x}\mathcal L v$, so $\mathcal A^*\phi=\lambda\phi\iff\mathcal L v=\lambda v$,
i.e.\ $v=u_0$ (same principal eigenfunction). Hence \eqref{eq:Fphi} and
$\phi F=e^{-\theta x}u_0\cdot e^{\theta x}u_0=u_0^2$.
\end{proof}

\begin{theorem}[General relative entropy / Lyapunov law]\label{thm:GRE}
Let $f(t,x)\ge0$ solve $\partial_t f=\mathcal A f$ with sufficient decay, and set the
\emph{renormalised density}
\begin{equation}\label{eq:renorm}
h(t,x):=\frac{f(t,x)\,e^{-\lambda t}}{F(x)} .
\end{equation}
Then $h$ solves the pure drift--diffusion equation
\begin{equation}\label{eq:hdrift}
\partial_t h=\frac{\widetilde D}{2}h'' + b(x)\,h',\qquad
b(x):=\widetilde D\frac{F'(x)}{F(x)}-\widetilde a=\widetilde D\frac{u_0'(x)}{u_0(x)} .
\end{equation}
For any convex $C^1$ function $H:[0,\infty)\to\R$, the \emph{generalised relative entropy}
\begin{equation}\label{eq:ent}
\mathcal H_H[t]:=\int_\R \phi(x)\,F(x)\,H\!\bigl(h(t,x)\bigr)\dd x
=\int_\R u_0^2\,H(h)\dd x
\end{equation}
is nonincreasing, with dissipation
\begin{equation}\label{eq:diss}
\frac{\mathrm d}{\mathrm dt}\mathcal H_H[t]
=-\frac{\widetilde D}{2}\int_\R u_0^2\,H''\!\bigl(h\bigr)\,|\partial_x h|^2\,\dd x\ \le\ 0 .
\end{equation}
\end{theorem}

\begin{proof}
Put $g:=e^{-\lambda t}f$, so $\partial_t g=(\mathcal A-\lambda)g$ and $h=g/F$. Then
$F\,\partial_t h=(\mathcal A-\lambda)(Fh)-h(\mathcal A-\lambda)F$. Since
$(\mathcal A-\lambda)F=0$ and a direct expansion gives
$(\mathcal A-\lambda)(Fh)-h(\mathcal A-\lambda)F=\frac{\widetilde D}{2}Fh''
+(\widetilde D F'-\widetilde a F)h'$, dividing by $F>0$ yields \eqref{eq:hdrift} with
$b=\widetilde D F'/F-\widetilde a$; and $F'/F=\theta+u_0'/u_0$ with $\widetilde D\theta=\widetilde a$
gives $b=\widetilde D u_0'/u_0$.

Differentiate \eqref{eq:ent}: since $\partial_t H(h)=H'(h)\partial_t h$,
\[
\frac{\mathrm d}{\mathrm dt}\mathcal H_H
=\int_\R u_0^2\,H'(h)\Bigl(\tfrac{\widetilde D}{2}h''+b h'\Bigr)\dd x .
\]
The weight $u_0^2$ is invariant for $\mathcal G:=\frac{\widetilde D}{2}\partial_x^2+b\partial_x$
in the sense that
\begin{equation}\label{eq:flux-id}
\partial_x\!\Bigl(\tfrac{\widetilde D}{2}u_0^2\,h'\Bigr)
=u_0^2\Bigl(\tfrac{\widetilde D}{2}h''+b h'\Bigr),
\end{equation}
which is verified directly: the left side is
$\frac{\widetilde D}{2}u_0^2h''+\widetilde D u_0u_0'h'
=u_0^2(\frac{\widetilde D}{2}h''+\widetilde D\frac{u_0'}{u_0}h')
=u_0^2(\frac{\widetilde D}{2}h''+b h')$. Therefore, integrating by parts (boundary terms vanish
by the super-polynomial decay of $u_0^2$, Theorem~\ref{thm:principal}(iv)),
\[
\frac{\mathrm d}{\mathrm dt}\mathcal H_H
=\int_\R H'(h)\,\partial_x\!\Bigl(\tfrac{\widetilde D}{2}u_0^2\,h'\Bigr)\dd x
=-\int_\R \tfrac{\widetilde D}{2}u_0^2\,H''(h)\,(h')^2\dd x\le0,
\]
since $H''\ge0$, $\widetilde D>0$, $u_0^2\ge0$. This is \eqref{eq:diss}.
\end{proof}

\begin{corollary}[Conservation and exponential convergence]\label{cor:conv}
With $H(s)=s$ in \eqref{eq:diss} the dissipation vanishes, giving
\begin{equation}\label{eq:cons-mass}
\frac{\mathrm d}{\mathrm dt}\int_\R \phi\, f\,e^{-\lambda t}\dd x=0,\qquad
\int_\R \phi\, f(t,\cdot)\dd x=e^{\lambda t}\int_\R\phi f_0 .
\end{equation}
Let $\gamma:=\lambda-\lambda_1>0$ be the spectral gap of $\mathcal L$, where $\lambda_1$ is its
second-largest eigenvalue (Lemma~\ref{lem:discrete}). Normalise $\int\phi f_0=\int\phi F=1$.
Then with $H(s)=(s-1)^2$,
\begin{equation}\label{eq:expconv}
\int_\R u_0^2\,(h-1)^2\,\dd x\ \le\ e^{-2\gamma t}\int_\R u_0^2\,(h_0-1)^2\,\dd x,
\end{equation}
so $e^{-\lambda t}f(t,\cdot)=Fh\to F$ exponentially fast (rate $\gamma$) in the weighted norm
$\|G\|^2:=\int (\phi/F)\,G^2=\int u_0^2(G/F)^2$.
\end{corollary}

\begin{proof}
\eqref{eq:cons-mass}: $H(s)=s$ gives $H''=0$, so $\frac{\mathrm d}{\mathrm dt}\int u_0^2 h=0$,
and $u_0^2 h=\phi F h=\phi g=\phi f e^{-\lambda t}$.

\eqref{eq:expconv}: the operator $\mathcal G=\frac{\widetilde D}{2}\partial_x^2+b\partial_x$ in
\eqref{eq:hdrift} is self-adjoint and $\le0$ on $L^2(u_0^2\dd x)$ with Dirichlet form
$\langle-\mathcal G\psi,\psi\rangle_{u_0^2}=\frac{\widetilde D}{2}\int u_0^2|\psi'|^2$ (by
\eqref{eq:flux-id} and IBP). The \emph{ground-state transform} $\psi\mapsto u_0\psi$ satisfies
the intertwining identity
$u_0\,\mathcal G\psi=(\mathcal L-\lambda)(u_0\psi)$, verified directly: with $V=\lambda-\frac{\widetilde D}{2}u_0''/u_0$,
$(\mathcal L-\lambda)(u_0\psi)=\frac{\widetilde D}{2}(u_0\psi)''+(V-\lambda)u_0\psi
=\frac{\widetilde D}{2}(u_0''\psi+2u_0'\psi'+u_0\psi'')-\frac{\widetilde D}{2}u_0''\psi
=u_0(\frac{\widetilde D}{2}\psi''+\widetilde D\frac{u_0'}{u_0}\psi')=u_0\mathcal G\psi$.
Hence $\mathcal G$ on $L^2(u_0^2\dd x)$ is unitarily equivalent (via $\psi\mapsto u_0\psi$, an
isometry $L^2(u_0^2\dd x)\to L^2(\dd x)$) to $\mathcal L-\lambda$ on $L^2(\dd x)$. The top
eigenvalue $0$ of $\mathcal G$ (eigenfunction $\psi\equiv1\mapsto u_0$, the ground state of
$\mathcal L$) is simple with spectral gap exactly $\lambda-\lambda_1=\gamma$ to the next
eigenvalue. Writing $w:=h-1$ with $\int u_0^2 w=0$ (by \eqref{eq:cons-mass}), the spectral gap
is the Poincar\'e inequality
$\frac{\widetilde D}{2}\int u_0^2|w'|^2\ge\gamma\int u_0^2 w^2$. Combined with
$\frac{\mathrm d}{\mathrm dt}\int u_0^2 w^2=-\widetilde D\int u_0^2|w'|^2$
(Theorem~\ref{thm:GRE}, $H(s)=(s-1)^2$, $H''=2$), this gives
$\frac{\mathrm d}{\mathrm dt}\int u_0^2 w^2\le-2\gamma\int u_0^2 w^2$ and \eqref{eq:expconv} by
Gr\"onwall. Finally $\int u_0^2(h-1)^2=\int(\phi/F)(Fh-F)^2=\|Fh-F\|^2\to0$.
\end{proof}

\section{(e) The inverse problem: recovering $\beta$ (and $\mu$) from $(\lambda,F)$}

We invert the forward map $(\beta,\mu)\mapsto(\lambda,F)$ from Theorem~\ref{thm:principal},
using the reduced model \eqref{eq:ode}.

\begin{theorem}[Inverse formula and identifiability]\label{thm:inverse}
Let data $(\lambda,F)$ be given with $F\in C^2(\R)$, $F>0$, $\int F=1$, $F\to0$ at $\infty$. Then:
\begin{enumerate}
\item[\rm(I1)] \textbf{(Unique explicit recovery of $m$.)}
\begin{equation}\label{eq:invm}
\boxed{\;m(x)=\beta(x)-\mu(x)=\lambda-\frac{\widetilde D}{2}\frac{F''(x)}{F(x)}+\widetilde a\,\frac{F'(x)}{F(x)}\;}
\qquad(F>0),
\end{equation}
the unique $m$ consistent with $(\lambda,F)$.
\item[\rm(I2)] \textbf{(Non-identifiability of $\beta,\mu$ separately.)} The data $(\lambda,F)$
determine \emph{only} $m=\beta-\mu$: for any smooth $\psi\ge0$, $(\beta+\psi,\mu+\psi)$ yields
the same $(\lambda,F)$. Hence $\beta$ (resp.\ $\mu$) is recoverable \emph{iff} $\mu$ (resp.\ $\beta$)
is known a priori, via $\beta=m+\mu$ (resp.\ $\mu=\beta-m$).
\item[\rm(I3)] \textbf{(Conditional stability.)} If $(\lambda_i,F_i)$, $i=1,2$, arise from $m_i$
with $F_i\ge c>0$ and $\|F_i\|_{C^2(K)}\le M$ on a compact $K$, then
\[
\|m_1-m_2\|_{L^\infty(K)}\le |\lambda_1-\lambda_2|
+\tfrac{\widetilde D}{2}\bigl\|\tfrac{F_1''}{F_1}-\tfrac{F_2''}{F_2}\bigr\|_{L^\infty(K)}
+|\widetilde a|\bigl\|\tfrac{F_1'}{F_1}-\tfrac{F_2'}{F_2}\bigr\|_{L^\infty(K)}
\le C(c,M)\bigl(|\lambda_1-\lambda_2|+\|F_1-F_2\|_{C^2(K)}\bigr).
\]
\end{enumerate}
\end{theorem}

\begin{proof}
\emph{(I1)} Since $F>0$, divide $\frac{\widetilde D}{2}F''-\widetilde a F'+mF=\lambda F$ by $F$
and solve for $m$; this gives \eqref{eq:invm} pointwise, uniquely. Smoothness of the right side
follows from $F\in C^\infty$ (Theorem~\ref{thm:principal}(iii)) and $F>0$.

\emph{(I2)} Adding $\psi$ to both $\beta,\mu$ leaves $m=\beta-\mu$ unchanged, hence (by (I1)) the
operator $\mathcal A$ and its principal pair $(\lambda,F)$. Conversely, if $(\beta_1,\mu_1)$ and
$(\beta_2,\mu_2)$ give the same $(\lambda,F)$, then \eqref{eq:invm} forces
$\beta_1-\mu_1=\beta_2-\mu_2$, i.e.\ they differ by $\psi:=\beta_1-\beta_2=\mu_1-\mu_2$. This
gauge ambiguity is the complete kernel of the inverse map.

\emph{(I3)} Subtract the two formulas \eqref{eq:invm} and take $L^\infty(K)$ norms (triangle
inequality) for the first bound. The maps $G\mapsto G''/G$, $G\mapsto G'/G$ are Lipschitz on
$\{G\in C^2(K):G\ge c,\ \|G\|_{C^2}\le M\}$: e.g.
$|\frac{F_1''}{F_1}-\frac{F_2''}{F_2}|\le\frac{|F_1''-F_2''|}{c}+\frac{M|F_1-F_2|}{c^2}$, and
similarly for the first-derivative term. Hence the right side is
$\le C(c,M)(|\lambda_1-\lambda_2|+\|F_1-F_2\|_{C^2(K)})$ with
$C(c,M)=\max\{1,(\tfrac{\widetilde D}{2}+|\widetilde a|)(\tfrac1c+\tfrac{M}{c^2})\}$.
\end{proof}

\begin{remark}[Resolving non-identifiability]
By (I2), steady-state data fix only $m=\beta-\mu$; full recovery needs an additional
independent measurement. Two natural options: (i) prior knowledge of one coefficient (e.g.\ a
measured $\mu$), after which \eqref{eq:invm} gives $\beta=m+\mu$; (ii) pre-limit data: in the
original model \eqref{eq:master}, $\beta$ multiplies the gain term $2\int\beta\varphi f$ while
$\mu$ enters only the loss term, so observing the population at distinct transition scales
$\varepsilon$ (or its second-moment growth) separates $\beta$ and $\mu$. The diffusion limit
\eqref{eq:ode} retains precisely the gauge-invariant datum $m=\beta-\mu$ and discards the rest.
\end{remark}

\section*{Numerical sanity check (bounded domain, Remark~\ref{rem:bounded})}

For $\widetilde D=1$, $\widetilde a=0.3$, $m(x)=1-\tfrac12 x^2$ on $(-12,12)$ with Dirichlet
data, a finite-difference discretisation of $\mathcal A$ yields a top eigenvalue
$\lambda\approx0.455$ whose eigenvector is strictly of one sign (positivity), with a spectral
gap $\approx1.0>0$, consistent with Theorem~\ref{thm:principal} and Corollary~\ref{cor:conv}.

\end{document}
